\documentclass[main.tex]{subfiles}

\begin{document}

\section{Conclusion\label{sec:conclusion}}

Interchange fees shape retail activity and redistribute consumption.
Using merchant-level data, we document substantial heterogeneity in the composition and cost of payments across stores, with critical implications for merchants' costs and redistribution.
By combining the results of a sufficient-statistics framework with a calibrated structural model, we quantify the overall impact of interchange pricing.
Specifically, we find that interchange fees redistribute approximately \$\absinput{combined_dollar_headline} billion per year from cash and debit users to credit card users, with consumer sorting attenuating these transfers by about \absinput{pct_reduction} percent.

Our analysis of policy counterfactuals reveals that recent changes in the payment system have had significant regressive impacts.
The Durbin Amendment, despite being designed to help consumers, primarily benefited credit card users rather than the intended debit card users, creating regressive wealth transfers from middle-income to high-income households.
Similarly, the rise of premium credit cards has been regressive, as higher-income consumers disproportionately capture the benefits.
These trends illustrate how network competition through richer rewards funded by higher fees amplifies transfers from low- to high-income households.
More broadly, these results highlight how interchange policies affect not only banks and card networks but also prices, market competition, and the distribution of economic surplus across consumers.

\end{document}